\chapter{Λοιπά σχόλια}

\section{Stack επιλογές και trade-offs}

Επιλέξαμε ένα \textbf{modern stack}: Python 3.13 / FastAPI / asyncmy στο
back-end και React 19 / TypeScript / Vite / TanStack στο front-end. Κάθε
συστατικό δικαιολογείται παρακάτω.

\begin{itemize}
    \item \textbf{Python 3.13 + asyncmy}. Επιλέξαμε \texttt{asyncmy} αντί
          \texttt{aiomysql} γιατί ο πρώτος έχει 4$\times$ throughput με
          identical API. Ανεξάρτητα από το αν ένα ORM θα ήταν ``πιο εύκολο''
          --- η εκφώνηση ζητάει \emph{raw SQL}, οπότε δεν τίθεται καν θέμα
          επιλογής. Επιπλέον, η επιλογή Python για το back-end
          ευθυγραμμίζεται με τη γλώσσα του ETL pipeline (Polars):
          αποφεύγουμε την εισαγωγή τρίτης γλώσσας στο project και
          κρατάμε ομογενές tooling, dependencies και testing
          infrastructure.
    \item \textbf{FastAPI + Pydantic v2}. Επιλέξαμε FastAPI για:
          (α) αυτόματη OpenAPI generation που τροφοδοτεί τους τύπους του
          front-end, (β) native async support, (γ) ξεκάθαρη separation
          μεταξύ schema validation και endpoint logic.
    \item \textbf{React 19 + TypeScript + Vite + TanStack}. Πιο μοντέρνο
          stack από Angular/Vue που είδαμε στο μάθημα, αλλά απολύτως
          κατάλληλο για single-page analytical UI με πολλά cached
          server-state queries. Το TanStack Query χειρίζεται caching,
          stale-while-revalidate, retry logic και background refetch.
          Ολο το front-end γράφτηκε σε \textbf{TypeScript}: typed DTOs
          που αντικατοπτρίζουν τα Pydantic schemas, typed search-params
          (Zod schemas σε κάθε route), strict null-checking. Ετσι τα
          contract breaks πιάνονται στο compile time αντί στο runtime.
    \item \textbf{Visx} αντί Chart.js / Recharts. Το Visx εκθέτει D3
          primitives ως React components, οπότε χτίσαμε custom chart
          components με ακριβή έλεγχο σε axis layout, tooltip, hover,
          drag-to-zoom, clip paths.
\end{itemize}

\section{Performance σημειώσεις}

Το \texttt{view\_author\_profile} αρχικά γύριζε σε περίπου 3 sec γιατί
έκανε \texttt{LEFT JOIN} στους \(\sim\)1.1M journal articles
\texttt{UNION ALL} στους \(\sim\)1.4M conference articles
\emph{για κάθε query}. Το λύσαμε με \texttt{materialized\_author\_profile}
που γεμίζει μία φορά μετά το ETL και έχει indexes σε
(\texttt{total\_articles}, \texttt{author\_name}). Τα authors browse
queries τώρα τρέχουν σε \(<\)30 ms.

Παρόμοιο pattern για \texttt{materialized\_authors\_vs\_articles\_scatter\_*}:
ένα DAO query αντί χιλιάδες point-by-point joins ανά request.

\section{Test coverage}

\begin{itemize}
    \item \textbf{Backend}: 200 pytest tests, 95\% line coverage. Πλήρης
          integration coverage εναντίον πραγματικού MySQL container
          (μέσω \texttt{conftest.py} που recreate το \texttt{mye030\_test}
          σε κάθε session).
    \item \textbf{Frontend unit / integration}: 160 vitest tests με
          MSW mocking. 39 test files. Καλύπτουν components, hooks,
          API client, routes και σενάρια filtering.
    \item \textbf{Frontend E2E}: ~50 Playwright tests σε 3 spec files
          (\texttt{app.spec.ts}, \texttt{charts.spec.ts},
          \texttt{report\_screenshots.spec.ts}) εναντίον live dev
          server. Καλύπτουν τα chart variants, τις 13 routes,
          σύστημα filters, sort/search interactions, zoom drag και reset,
          legend toggle, checkbox multi-select. Το
          \texttt{report\_screenshots.spec.ts} παράγει αυτόματα τα PNGs
          που χρησιμοποιούνται στο Κεφάλαιο 4.
\end{itemize}

\section{Στρατηγική joins και indexes}

Επειδή κάθε analytical query του front-end ενώνει fact με lookup ---
και συχνά περνά μέσα από bridge --- οι αποφάσεις γύρω από τα joins
και τα indexes είναι κρίσιμες.

\subsection{Join patterns}

\begin{itemize}\setlength{\itemsep}{0.2em}
    \item \textbf{Profile σελίδα} (journal/conference):
          \texttt{fact\_*\_articles INNER JOIN lookup\_*}. INNER γιατί
          το FK εγγυάται την ύπαρξη του parent venue.
    \item \textbf{Article authors list}:
          \texttt{fact\_*\_articles LEFT JOIN bridge\_*\_authors LEFT JOIN lookup\_authors}.
          LEFT για να εμφανίζονται και άρθρα χωρίς authors (orphans).
    \item \textbf{Year article list}:
          \texttt{fact\_journal\_articles UNION ALL fact\_conference\_articles}
          (φιλτραρισμένο ανά έτος), μετά INNER JOIN με τα venues.
          \texttt{UNION ALL} (όχι \texttt{UNION}) γιατί τα δύο fact
          tables έχουν επικαλυπτόμενα \texttt{article\_id} ranges,
          οπότε \texttt{UNION} θα έσβηνε γνήσια διακριτές γραμμές.
    \item \textbf{Author profile}:
          \texttt{lookup\_authors INNER JOIN materialized\_author\_profile}.
          Το union των δύο bridges γίνεται μία φορά στο ETL, ποτέ
          per request.
    \item \textbf{Charts (subject area, FoR, scatter)}: αναγνώσεις
          από \texttt{materialized\_*} αντί για on-the-fly aggregate
          των views.
\end{itemize}

\subsection{Επιλογές INNER vs LEFT}

Κάθε join διαλέχτηκε με βάση τις εγγυήσεις του FK, όχι αντανακλαστικά:

\begin{itemize}\setlength{\itemsep}{0.2em}
    \item Λάθος LEFT $\rightarrow$ φουσκωμένος row count, λάθος metrics.
    \item Λάθος INNER $\rightarrow$ σιωπηρή απώλεια legitimate rows.
\end{itemize}

\texttt{COUNT(DISTINCT)} εμφανίζεται μόνο μέσα στο materialised path,
ποτέ στο per-request path.

\subsection{Indexes πέρα από τα PK/FK}

\begin{center}
\begin{tabular}{@{}p{0.55\textwidth} p{0.40\textwidth}@{}}
\toprule
\textbf{Index} & \textbf{Καλύπτει} \\
\midrule
\texttt{fact\_journal\_articles (journal\_id, year)}
& Journal yearly stats, articles list, year-range φίλτρα \\
\texttt{fact\_conference\_articles (conference\_id, year)}
& Conference equivalents \\
\texttt{bridge\_*\_authors (author\_id, article\_id)}
& Author article list, ETL builds \\
\texttt{lookup\_journals (best\_subject\_area)},
\texttt{lookup\_journals (publisher)}
& Dropdown φίλτρα \\
\texttt{materialized\_author\_profile (total\_articles DESC, author\_name)}
& Authors browse default sort χωρίς filesort \\
\bottomrule
\end{tabular}
\end{center}

\texttt{EXPLAIN} δείχνει zero ``Using filesort'' για όλα τα default
sort orders. Οι writes (one-shot ETL load) πληρώνουν το κόστος των
indexes --- αποδεκτό για read-mostly workload.

\section{Πίνακας trade-offs}

Συνοπτική παρουσίαση όλων των κρίσιμων αποφάσεων αρχιτεκτονικής
με τα trade-offs τους. Πλήρης ανάλυση (alternatives considered,
context, consequences) στο \texttt{docs/ARCHITECTURE.md}.

\begin{center}
\begin{tabular}{@{}p{0.30\textwidth} p{0.32\textwidth} p{0.32\textwidth}@{}}
\toprule
\textbf{Απόφαση} & \textbf{Γιατί} & \textbf{Κόστος που δεχτήκαμε} \\
\midrule
Raw SQL μέσω asyncmy, χωρίς ORM
& Απαίτηση εκφώνησης + 4$\times$ throughput vs aiomysql
& Κάθε endpoint = χειρόγραφο DAO method \\
\addlinespace
Wide journal dimension
& One-row reads, απλά indexes
& $\sim$170 KB επιπλέον χώρος \\
\addlinespace
Materialised aggregates
& 11--46 s $\rightarrow$ < 0.3 s queries
& Read-only μεταξύ ETL runs \\
\addlinespace
Hybrid entity resolution
& 100\% precision σε acronyms, recall σε fuzzy titles
& 81\% conferences unranked (ορατό μέσω \texttt{ranked\_only}) \\
\addlinespace
FastAPI + Pydantic v2
& Δωρεάν OpenAPI, async, RFC 7807 errors
& Δύο type systems (Pydantic + TS) συντηρούνται με το χέρι \\
\addlinespace
Visx πάνω από Recharts
& Pixel-perfect charts με drag-to-zoom + Voronoi hover
& Περισσότερος component κώδικας ανά chart \\
\addlinespace
URL-driven filters (TanStack)
& Shareable URLs, browser back, no flicker
& Schema validation (Zod) σε κάθε route \\
\addlinespace
RFC 7807 problem-details
& Single error envelope, machine-readable
& Custom handler στο \texttt{api/errors.py} \\
\addlinespace
Polars αντί pandas
& $\sim$10$\times$ faster, lower memory
& Νεότερο ecosystem \\
\addlinespace
Static post-ETL warehouse
& Ταιριάζει με το read-mostly workload
& Δεν υπάρχει live ingest, το ETL είναι re-run \\
\bottomrule
\end{tabular}
\end{center}

\section{Ενιαίο orchestrator script: \texttt{python run.py}}

Το αρχείο \texttt{run.py} στο repo root είναι ένα stdlib-only
orchestrator (~250 γραμμές, Python 3.10+, χωρίς third-party deps)
που αυτο-εντοπίζει τι υπάρχει στον δίσκο και σηκώνει αυτόματα το
σωστό σενάριο:

\begin{itemize}\setlength{\itemsep}{0.2em}
    \item \textbf{Σενάριο A}: εάν υπάρχει
          \texttt{deliverables/db\_backup.sql.gz}, κάνει
          auto-restore στο first boot ($\sim$2.5 λεπτά).
    \item \textbf{Σενάριο B}: εάν υπάρχουν μόνο τα raw CSVs στο
          \texttt{data/}, τρέχει την ETL pipeline μέσα στο backend
          container και φορτώνει τα 3 SQL scripts ($\sim$5 λεπτά).
    \item \textbf{Σενάριο C}: εάν δεν υπάρχει τίποτα, εκτυπώνει
          banner με Drive link + DBLP/iCore/Kaggle URLs και κάνει
          exit. Καμία σιωπηρή αποτυχία.
\end{itemize}

Πρόσθετα entry points: \texttt{python run.py -{}-status} (health των
containers + smoke probes), \texttt{-{}-verify} (data quality
report), \texttt{-{}-etl} (force ETL), \texttt{-{}-down} (stop with
volume kept), \texttt{-{}-reset} (stop + wipe).

\section{Δύο τρόποι εκτέλεσης: production vs development}

Η εφαρμογή προσφέρεται με δύο παράλληλους τρόπους εκτέλεσης που
καλύπτουν διαφορετικές ανάγκες:

\begin{itemize}
    \item \textbf{Production stack (Path 0)} --- ένα και μόνο
          \texttt{docker compose up -d -{}-wait} (ή
          \texttt{python run.py}). Χτίζει τρία images
          (mye030\_mysql, mye030\_backend, mye030\_frontend) και
          τα σηκώνει με healthchecks και \texttt{depends\_on}
          ordering. Το frontend image είναι multi-stage build: το
          \texttt{vite build} παράγει minified, content-hashed bundle
          ($\sim$600 KB gzipped) που σερβίρει η nginx 1.27-alpine.
          Η nginx επιπλέον κάνει reverse-proxy τα
          \texttt{/api/*}, \texttt{/docs}, \texttt{/redoc},
          \texttt{/openapi.json} και \texttt{/health} στο backend,
          οπότε όλο το προϊόν ζει πίσω από ένα μόνο origin (χωρίς
          ανάγκη για CORS configuration).
    \item \textbf{Development stack (Paths A/B/C)} ---
          host-installed Python + Node, με \texttt{uvicorn -{}-reload}
          στο backend και \texttt{pnpm dev} (Vite dev server) στο
          front-end. Hot module replacement, source maps, on-demand
          transpilation. Κατάλληλο για ενεργή ανάπτυξη.
\end{itemize}

\textbf{Port conflict resilience.} Οι τρεις host ports
(\texttt{FRONTEND\_PORT=5173}, \texttt{BACKEND\_PORT=8000},
\texttt{MYSQL\_PORT=3306}) είναι παραμετροποιήσιμες μέσω
\texttt{.env}. Το \texttt{scripts/setup.ps1} / \texttt{setup.sh}
κάνει pre-flight έλεγχο και βγαίνει με ξεκάθαρο μήνυμα που
ονομάζει την διεργασία που μπλοκάρει. Σε περίπτωση clash, ο
χρήστης απλά αλλάζει την αντίστοιχη τιμή στο \texttt{.env}
(π.χ. \texttt{FRONTEND\_PORT=5273}) και ξανατρέχει
\texttt{docker compose up}.

\section{Παρατηρήσιμες πηγές (τι μπορεί να δει ο εξεταστής)}

Όλες οι αξιώσεις του παραπάνω πίνακα είναι ελέγξιμες:

\begin{itemize}\setlength{\itemsep}{0.2em}
    \item \textbf{Live app}: \texttt{http://localhost:5173}
    \item \textbf{Swagger UI}: \texttt{http://localhost:8000/docs}
    \item \textbf{ReDoc}: \texttt{http://localhost:8000/redoc}
    \item \textbf{OpenAPI JSON} (live ή offline στο \texttt{docs/openapi.json})
    \item \textbf{HTTP playbook}: \texttt{src/backend/api.http} με REST Client
    \item \textbf{Backend tests}: \texttt{uv run pytest tests/ -q}
    \item \textbf{Coverage HTML}: \texttt{uv run pytest -{}-cov -{}-cov-report=html}
    \item \textbf{Frontend tests}: \texttt{pnpm test} / \texttt{pnpm test:coverage}
    \item \textbf{E2E suite}: \texttt{pnpm exec playwright test}
    \item \textbf{E2E HTML report}: \texttt{pnpm exec playwright show-report}
    \item \textbf{Lint + types}: \texttt{uv run ruff check . \&\& uv run pyright}
    \item \textbf{Database CLI}: \texttt{docker exec -it mye030\_mysql mysql -uroot -proot mye030}
\end{itemize}

\section{Επόμενα βήματα}

\begin{itemize}
    \item Cmd+K command palette για global search (deferred).
    \item Project structure refactor σε
          \texttt{features/} domain-driven slices (deferred).
\end{itemize}

\section{Αναφορές}

\begin{itemize}
    \item DBLP. \url{https://dblp.org/}
    \item Kaggle Scimago journal ranking dataset.
    \item iCore26 conference rankings. \url{http://portal.core.edu.au/conf-ranks/}
    \item FastAPI documentation. \url{https://fastapi.tiangolo.com/}
    \item Visx. \url{https://airbnb.io/visx/}
    \item TanStack Router \& Query. \url{https://tanstack.com/}
    \item asyncmy driver. \url{https://github.com/long2ice/asyncmy}
\end{itemize}
